% Options for packages loaded elsewhere
\PassOptionsToPackage{unicode}{hyperref}
\PassOptionsToPackage{hyphens}{url}
\documentclass[
  ignorenonframetext,
]{beamer}
\newif\ifbibliography
\usepackage{pgfpages}
\setbeamertemplate{caption}[numbered]
\setbeamertemplate{caption label separator}{: }
\setbeamercolor{caption name}{fg=normal text.fg}
\beamertemplatenavigationsymbolsempty
% remove section numbering
\setbeamertemplate{part page}{
  \centering
  \begin{beamercolorbox}[sep=16pt,center]{part title}
    \usebeamerfont{part title}\insertpart\par
  \end{beamercolorbox}
}
\setbeamertemplate{section page}{
  \centering
  \begin{beamercolorbox}[sep=12pt,center]{section title}
    \usebeamerfont{section title}\insertsection\par
  \end{beamercolorbox}
}
\setbeamertemplate{subsection page}{
  \centering
  \begin{beamercolorbox}[sep=8pt,center]{subsection title}
    \usebeamerfont{subsection title}\insertsubsection\par
  \end{beamercolorbox}
}
% Prevent slide breaks in the middle of a paragraph
\widowpenalties 1 10000
\raggedbottom
\AtBeginPart{
  \frame{\partpage}
}
\AtBeginSection{
  \ifbibliography
  \else
    \frame{\sectionpage}
  \fi
}
\AtBeginSubsection{
  \frame{\subsectionpage}
}
\usepackage{iftex}
\ifPDFTeX
  \usepackage[T1]{fontenc}
  \usepackage[utf8]{inputenc}
  \usepackage{textcomp} % provide euro and other symbols
\else % if luatex or xetex
  \usepackage{unicode-math} % this also loads fontspec
  \defaultfontfeatures{Scale=MatchLowercase}
  \defaultfontfeatures[\rmfamily]{Ligatures=TeX,Scale=1}
\fi
\usepackage{lmodern}
\ifPDFTeX\else
  % xetex/luatex font selection
\fi
% Use upquote if available, for straight quotes in verbatim environments
\IfFileExists{upquote.sty}{\usepackage{upquote}}{}
\IfFileExists{microtype.sty}{% use microtype if available
  \usepackage[]{microtype}
  \UseMicrotypeSet[protrusion]{basicmath} % disable protrusion for tt fonts
}{}
\makeatletter
\@ifundefined{KOMAClassName}{% if non-KOMA class
  \IfFileExists{parskip.sty}{%
    \usepackage{parskip}
  }{% else
    \setlength{\parindent}{0pt}
    \setlength{\parskip}{6pt plus 2pt minus 1pt}}
}{% if KOMA class
  \KOMAoptions{parskip=half}}
\makeatother
\usepackage{longtable,booktabs,array}
\usepackage{caption}
\captionsetup[table]{skip=6pt}
\newcounter{none} % for unnumbered tables
\usepackage{calc} % for calculating minipage widths
\usepackage{caption}
% Make caption package work with longtable
\makeatletter
\def\fnum@table{\tablename~\thetable}
\makeatother
\setlength{\emergencystretch}{3em} % prevent overfull lines
\providecommand{\tightlist}{%
  \setlength{\itemsep}{0pt}\setlength{\parskip}{0pt}}
% Manuscript macro/package fallbacks (newcommand -> providecommand).
\usepackage{amsmath}
\usepackage{amssymb}
\usepackage{booktabs}
% Contents listing: article.cls prefixes every \l@section entry with
% \addvspace{1.0em}, which across ten sections costs about ten lines and pushed
% the ~40-entry listing one line past page 2 — leaving a third sheet carrying a
% single "10 References" line. Tighten that inter-section gap for the duration
% of \tableofcontents only, inside a group, so body spacing is untouched and no
% entry is dropped (reducing \tocdepth would have hidden 29 real subsections).
  \providecommand{\addvspace}[1]{\vskip 2\p@}%
% Document metadata. config.yaml declares a subtitle and a keyword list, and
% the producer emits both as tokens, but the template's \hypersetup writes only
% pdftitle/pdfauthor/pdflang -- so `pdfinfo` reported no Subject and no
% Keywords. This block runs after the template's, so it wins.
%
% The two values below are WRITTEN by scripts/z_generate_manuscript_variables.py
% (sync_preamble_metadata), never typed. A double-brace token cannot be used here:
% preamble.md is substituted into output/manuscript/ like any other section, but
% the template's _manuscript_source.py then copies the raw file back over that
% copy, so an unresolved token would reach hyperref and land in the PDF's
% metadata verbatim.
% Bounded identifier splitting. The template defines
%   \protected\def\breaktt#1{\begingroup\ttfamily\seqsplit{#1}\endgroup}
% and \seqsplit permits a break after EVERY character with no continuation cue.
% In the 2026-09-05 PDF that rendered `model_family_manifest.json` as
% `model_family_manifest.js` + `on` and `src/manuscript_variables.py` as
% `(src/manusc` + `ript_variables.py` — the first is itself a valid filename, so
% a reader cannot tell a line break from a name. Keep seqsplit's guarantee that
% no code token overflows the measure, but forbid breaks inside the last
% \GNNttTail characters (an extension is never orphaned) and inside the first
% \GNNttHead (a one- or two-letter stub is never left behind).

% Slide layout defaults for warning-clean scientific decks.
\usepackage{etoolbox}
\IfFileExists{xurl.sty}{\usepackage{xurl}}{}
\IfFileExists{seqsplit.sty}{\usepackage{seqsplit}}{\newcommand{\seqsplit}[1]{#1}}
\protected\def\breakseq#1{\seqsplit{#1}}
\protected\def\breaktt#1{\begingroup\ttfamily\seqsplit{#1}\endgroup}
\setlength{\emergencystretch}{6em}
\tolerance=5000
\hbadness=10000
\hfuzz=1pt
\setlength{\tabcolsep}{2pt}
\AtBeginEnvironment{longtable}{\tiny\renewcommand{\arraystretch}{0.86}\setlength{\tabcolsep}{1pt}}
\AtBeginEnvironment{tabular}{\tiny\renewcommand{\arraystretch}{0.86}\setlength{\tabcolsep}{1pt}}
\AtBeginEnvironment{equation}{\tiny}
\AtBeginEnvironment{equation*}{\tiny}
\AtBeginEnvironment{align}{\tiny}
\AtBeginEnvironment{align*}{\tiny}
\AtBeginEnvironment{itemize}{\footnotesize}
\AtBeginEnvironment{enumerate}{\footnotesize}
\AtBeginEnvironment{description}{\footnotesize}
\setbeamerfont{caption}{size=\tiny}
\setbeamerfont{caption name}{size=\tiny}
\setbeamerfont{normal text}{size=\small}
\setbeamerfont{frametitle}{size=\small}
\setbeamerfont{section title}{size=\footnotesize}
\setbeamerfont{subsection title}{size=\footnotesize}
\setbeamertemplate{section page}{%
  \centering
  \begin{beamercolorbox}[sep=12pt,center,wd=\paperwidth]{section title}
    \parbox{0.86\paperwidth}{\centering\usebeamerfont{section title}\insertsection\par}
  \end{beamercolorbox}
}
\setbeamertemplate{subsection page}{%
  \centering
  \begin{beamercolorbox}[sep=8pt,center,wd=\paperwidth]{subsection title}
    \parbox{0.86\paperwidth}{\centering\usebeamerfont{subsection title}\insertsubsection\par}
  \end{beamercolorbox}
}
\setlength{\abovecaptionskip}{2pt}
\setlength{\belowcaptionskip}{0pt}

% Natbib and cross-reference fallbacks — slides don't load natbib
% or cleveref, but combined-PDF manuscript prose may emit these
% commands. The fallback renders citations as a bracketed key list
% and cross-references as detokenized labels so slides don't fail on
% undefined control sequences or raw underscores. \providecommand is
% a no-op if packages load later.
\providecommand{\citep}[1]{[#1]}
\providecommand{\citet}[1]{#1}
\providecommand{\citealp}[1]{#1}
\providecommand{\citeauthor}[1]{#1}
\providecommand{\citeyear}[1]{#1}
\providecommand{\cref}[1]{\texttt{\detokenize{#1}}}
\providecommand{\Cref}[1]{\texttt{\detokenize{#1}}}
\providecommand{\crefrange}[2]{\texttt{\detokenize{#1}}--\texttt{\detokenize{#2}}}
\providecommand{\Crefrange}[2]{\texttt{\detokenize{#1}}--\texttt{\detokenize{#2}}}
\providecommand{\paragraph}[1]{\textbf{#1}\ }

\newtheorem{proposition}{Proposition}
\newtheorem{hypothesis}{Hypothesis}
\newtheorem{remark}[theorem]{Remark}
\newtheorem{axiom}{Axiom}
\newtheorem{property}{Property}
\usepackage{bookmark}
\IfFileExists{xurl.sty}{\usepackage{xurl}}{} % add URL line breaks if available
\urlstyle{same}
% fallback for those not using the hyperref driver hyperxmp:
\makeatletter
\@ifundefined{xmpquote}{\newcommand{\xmpquote}[1]{#1}}{}
\makeatother
\hypersetup{
  hidelinks,
  pdfcreator={LaTeX via pandoc}}

\author{\texorpdfstring{}{}}
\date{}

\begin{document}

\section{Symbols and Glossary}\label{sec:symbols_glossary}

\begin{frame}[allowframebreaks]{Symbols and Glossary}
This glossary defines the Generalized Notation Notation (GNN) language
constructs used to specify generative models, together with the Active
Inference symbols those constructs carry across the model kinds of
sec.~3.2. Definitions follow the GNN syntax
specification and the exemplar specifications distributed with the
repository --- the discrete POMDP exemplars and the continuous
linear-Gaussian exemplars alike --- and they align with the standard
formulations of both kinds (Smékal and Friedman 2023; Da Costa et al.
2020; Parr et al. 2022).
\end{frame}

\begin{frame}[fragile,allowframebreaks]{Language Constructs}
\protect\phantomsection\label{language-constructs}
A GNN file is an ordered, UTF-8 Markdown document whose level-2 headers
name required and optional sections. The strict schema validator
enforces section order, declaration grammar, and connection syntax; the
constructs below are the load-bearing sections a parser must recognize;
they are catalogued in tbl.~\ref{tbl:gnn_constructs}.

\end{frame}

\begin{frame}[fragile,allowframebreaks]{Language Constructs}
\begin{longtable}[]{@{}
  >{\raggedright\arraybackslash}p{(\linewidth - 2\tabcolsep) * \real{0.5000}}
  >{\raggedright\arraybackslash}p{(\linewidth - 2\tabcolsep) * \real{0.5000}}@{}}
\caption{The GNN language constructs a conforming parser must recognize,
with the meaning each carries. The first rows are the required sections
that pin a model's identity, variables, and factor graph; the optional
sections that follow carry parameterization, ontology bindings, timing,
equations, and provenance. The operator rows are the connection
vocabulary the \protect\texttt{Connections} section is written in ---
directed \protect\texttt{A\textgreater{}B}, undirected
\protect\texttt{A-B}, and their labeled v1.1 forms --- and
\protect\texttt{default=…} supplies initialization hints. Definitions
follow the GNN syntax specification and the exemplar specifications
under \protect\texttt{input/gnn\_files}, and the strict schema validator
enforces section order and declaration grammar before any downstream
step runs. The construct vocabulary is deliberately kind-agnostic: the
same \protect\texttt{StateSpaceBlock}, \protect\texttt{Connections}, and
operator syntax declare a discrete categorical model, a continuous
linear-Gaussian model, or a composed multi-agent or hierarchical one ---
what differs between kinds is which matrix names the block declares, as
tbl.~\ref{tbl:actinf_symbols}
details.}\label{tbl:gnn_constructs}\tabularnewline
\toprule\noalign{}
\begin{minipage}[b]{\linewidth}\raggedright
Construct
\end{minipage} & \begin{minipage}[b]{\linewidth}\raggedright
Meaning
\end{minipage} \\
\midrule\noalign{}
\endfirsthead
\toprule\noalign{}
\begin{minipage}[b]{\linewidth}\raggedright
Construct
\end{minipage} & \begin{minipage}[b]{\linewidth}\raggedright
Meaning
\end{minipage} \\
\midrule\noalign{}
\endhead
\texttt{\#\#\ GNNSection} & Required short identifier for the model
block (no spaces, e.g.~\texttt{ActInfPOMDP}). \\
\breaktt{\#\#\ GNNVersionAndFlags} & Required version declaration
(\texttt{GNN\ v1} or \texttt{GNN\ v1.1}) with optional flags. \\
\texttt{\#\#\ ModelName} & Required human-readable title of the
generative model. \\
\breaktt{\#\#\ StateSpaceBlock} & Required block declaring every variable
and matrix as \texttt{NAME{[}dim,\ …,\ type=…{]}}, one per line --- the
block whose declared vocabulary fixes the model's kind. \\
\texttt{\#\#\ Connections} & Required edge list relating state-space
variables, expressing the model's factor graph. \\
\breaktt{\#\#\ ModelAnnotation} & Optional free-text description of the
model's modalities, factors, and assumptions. \\
\breaktt{\#\#\ InitialParameterization} & Optional concrete numeric
values for declared matrices and vectors. \\
\breaktt{\#\#\ ActInfOntologyAnnotation} & Optional bindings from each
variable to a CamelCase Active Inference ontology term. \\
\breaktt{\#\#\ ModelParameters} & Optional key-value dimensions
(e.g.~\breaktt{num\_hidden\_states}, \texttt{num\_obs},
\texttt{num\_actions}) consumed by code generators. \\
\texttt{\#\#\ Time} & Optional dynamics declaration: a time variable
plus \texttt{Dynamic}/\texttt{Static},
\texttt{Discrete}/\texttt{Continuous}, and \breaktt{ModelTimeHorizon}. \\
\texttt{\#\#\ Equations} & Optional LaTeX-rendered formulas defining
model dynamics and the relationships between declared variables;
round-tripped by the semantic-fidelity gate. \\
\texttt{\#\#\ Footer} & Optional closing section that terminates the
file and lets a reader or parser enter it from either end. \\
\texttt{\#\#\ Signature} & Optional provenance block carrying a
cryptographic signature over the specification. \\
\texttt{NAME{[}d₁,d₂,…,type=…{]}} & A variable or tensor declaration;
dimensions are positive integers or named references, and a
\texttt{type} (\texttt{float}, \texttt{int}, \texttt{bool}) is
required. \\
\texttt{A\textgreater{}B} & Directed (causal) connection operator: edge
from \texttt{A} to \texttt{B}, e.g.~\texttt{D\textgreater{}s} (a prior
conditions a hidden state) or \texttt{F\textgreater{}x} (a transition
matrix drives a continuous state). \\
\texttt{A-B} & Undirected (bidirectional) connection operator,
e.g.~\texttt{s-A} (a hidden state participates in the likelihood
mapping). \\
\texttt{A\textgreater{}B:label} / \texttt{A-B:label} & A v1.1 annotated
edge; the trailing label documents the relation and is preserved but may
be ignored for structural validation. \\
\texttt{default=…} & A v1.1 declaration hint (\texttt{uniform},
\texttt{zeros}, \texttt{ones}, \texttt{eye}, \texttt{random}) supplying
an initialization for a matrix or vector. \\
\bottomrule\noalign{}
\end{longtable}
\end{frame}

\begin{frame}[fragile,allowframebreaks]{Symbols by Model Kind}
\protect\phantomsection\label{symbols-by-model-kind}
The exemplar specifications declare the generative-model components of
Active Inference, mapping each GNN variable to its probabilistic meaning
(Da Costa et al. 2020; Smith et al. 2022; Parr et al. 2022). Because the
notation spans model kinds, the same letter can name different objects
in different kinds --- most prominently \texttt{F}, which is the
variational free energy of eq.~6 in the discrete kind and the
state-transition matrix of eq.~9 in the continuous
kind --- so each row of tbl.~\ref{tbl:actinf_symbols} names the kind its
reading belongs to and the equation in sec.~3
that fixes it, and the glossary and the formal statement cannot drift
apart.

\end{frame}

\begin{frame}[fragile,allowframebreaks]{Symbols by Model Kind}
\begin{longtable}[]{@{}
  >{\raggedright\arraybackslash}p{(\linewidth - 2\tabcolsep) * \real{0.5000}}
  >{\raggedright\arraybackslash}p{(\linewidth - 2\tabcolsep) * \real{0.5000}}@{}}
\caption{The Active Inference symbols a GNN specification carries,
scoped by model kind and bound to the equation of
sec.~3 that fixes each role. Read the
discrete-kind rows as the categorical generative-model tensors
(likelihood, controlled transitions, log-preferences, prior over initial
states, habit prior) with \protect\texttt{s}, \protect\texttt{o},
\protect\texttt{\texorpdfstring{\ensuremath{\pi}}{pi}}, and \protect\texttt{u} as the inference variables
and \protect\texttt{F} and \protect\texttt{G} as the two free-energy
functionals minimized at eq.~6 and eq.~7; read the
continuous-kind rows as the linear-Gaussian parameterization --- system
matrices, covariances, Gaussian prior, and the optional closed-loop
declarations of eq.~11. The rows where one letter
carries two kind-scoped readings (\protect\texttt{F},
\protect\texttt{u}, \protect\texttt{t}) are deliberate: the notation
reuses a compact alphabet across kinds, and the declared matrix
vocabulary of the \protect\texttt{StateSpaceBlock} --- not the letter
alone --- is what fixes the reading, which is exactly what the kind
classification of sec.~3.2 computes. Use the table as
a decoding key when reading a \protect\texttt{StateSpaceBlock}: every
declared matrix name should map to exactly one kind-scoped row here, and
the \protect\breaktt{\#\#\ ActInfOntologyAnnotation} bindings described
below are what make that mapping machine-checkable
downstream.}\label{tbl:actinf_symbols}\tabularnewline
\toprule\noalign{}
\begin{minipage}[b]{\linewidth}\raggedright
Symbol
\end{minipage} & \begin{minipage}[b]{\linewidth}\raggedright
Meaning
\end{minipage} \\
\midrule\noalign{}
\endfirsthead
\toprule\noalign{}
\begin{minipage}[b]{\linewidth}\raggedright
Symbol
\end{minipage} & \begin{minipage}[b]{\linewidth}\raggedright
Meaning
\end{minipage} \\
\midrule\noalign{}
\endhead
\texttt{A} & Discrete kind: likelihood (observation) matrix encoding
\(P(o \mid s)\), a column-stochastic map from hidden states to
observation outcomes (eq.~2). \\
\texttt{B} & Discrete kind: transition tensor encoding
\(P(s' \mid s, u)\), one column-stochastic slice per action in the
canonical \texttt{B{[}next\_state{]}{[}previous\_state{]}{[}action{]}}
order (eq.~3). \\
\texttt{C} & Discrete kind: preference vector --- log-preferences over
observation outcomes that bias the agent toward preferred outcomes
(eq.~4). \\
\texttt{D} & Discrete kind: prior vector over initial hidden states
(eq.~5). \\
\texttt{E} & Discrete kind: habit vector --- an initial policy prior
over actions, entering the policy posterior (eq.~8). \\
\texttt{s} & Discrete kind: current hidden-state distribution;
\texttt{s\_prime} (\texttt{s\textquotesingle{}}) is the next
hidden-state distribution. \\
\texttt{o} & Discrete kind: current observation, an integer index over
outcome modalities. \\
\texttt{\texorpdfstring{\ensuremath{\pi}}{pi}} & Discrete kind: policy --- a distribution over actions
inferred from expected free energy (eq.~8). \\
\texttt{u} & Discrete kind: the selected (sampled) action, sampled from
the policy posterior. Continuous kind: the control input added to the
state each step (eq.~9). \\
\texttt{F} & Discrete kind: variational free energy, minimized during
state inference (eq.~6) (Friston 2010). Continuous kind: the
state-transition matrix of the linear-Gaussian model
(eq.~9). \\
\texttt{G} & Discrete kind: expected free energy per policy, minimized
during policy inference to score candidate actions (eq.~7)
(Da Costa et al. 2020). \\
\texttt{t} & Discrete time step; the horizon \(T\) bounds the product in
eq.~1. Continuous kind: the time index of the
state-space model. \\
\texttt{H} & Continuous kind: observation matrix mapping the latent
state to the observation (eq.~10). \\
\texttt{Q} & Continuous kind: process-noise covariance of the state
evolution (eq.~9). \\
\texttt{R} & Continuous kind: observation-noise covariance of the
readout (eq.~10). \\
\texttt{x} & Continuous kind: the Gaussian latent state \(x_t\);
\texttt{x{[}2,1,type=float{]}} in the navigation exemplar declares a
two-dimensional position. \\
\texttt{y} & Continuous kind: the continuous observation read out
linearly from the state (eq.~10). \\
\texttt{prior\_mean} & Continuous kind: prior mean \(\mu_0\) over the
initial latent state (eq.~9). \\
\texttt{prior\_cov} & Continuous kind: prior covariance \(\Sigma_0\)
over the initial latent state (eq.~9). \\
\texttt{goal\_mean} & Continuous kind: preferred state \(\mu^\star\) the
closed-loop controller steers toward (eq.~11). \\
\texttt{control\_gain} & Continuous kind: scalar proportional gain \(k\)
on the belief error in the closed loop (eq.~11). \\
\bottomrule\noalign{}
\end{longtable}
\end{frame}

\begin{frame}[fragile,allowframebreaks]{Ontology Bindings and
Implementations}
\protect\phantomsection\label{ontology-bindings-and-implementations}
The \breaktt{\#\#\ ActInfOntologyAnnotation} section binds each variable
to a canonical term --- \breaktt{A=LikelihoodMatrix},
\breaktt{B=TransitionMatrix}, \breaktt{C=LogPreferenceVector},
\breaktt{D=PriorOverHiddenStates}, \texttt{s=HiddenState},
\texttt{o=Observation}, \texttt{\texorpdfstring{\ensuremath{\pi}}{pi}=PolicyVector} in the discrete
exemplars; \breaktt{F=StateTransitionMatrix},
\breaktt{H=ObservationMatrix}, \breaktt{Q=ProcessNoiseCovariance},
\breaktt{R=ObservationNoiseCovariance}, \breaktt{prior\_mean=PriorMean},
\breaktt{prior\_cov=PriorCovariance}, \breaktt{goal\_mean=PreferredState},
\breaktt{control\_gain=ControlGain}, \breaktt{x=ContinuousHiddenState},
\breaktt{y=ContinuousObservation}, \texttt{u=ControlInput} in the
continuous exemplars --- which downstream pipeline steps use for
semantic analysis and validation. The bindings are per-kind by
\end{frame}

\begin{frame}[fragile,allowframebreaks]{Ontology Bindings and
Implementations}
construction: a continuous exemplar binds \texttt{F} to
\breaktt{StateTransitionMatrix}, never to a free-energy term, so the
ontology layer carries the same kind-scoped reading
tbl.~\ref{tbl:actinf_symbols} fixes by hand. The same GNN specification
feeds the project's rendering backends --- including the 9 registered
targets (PyMDP, RxInfer.jl, ActiveInference.jl, JAX, DisCoPy, PyTorch,
NumPyro, Stan, bnlearn), of which the 9 listed as executable (PyMDP,
RxInfer.jl, ActiveInference.jl, JAX, DisCoPy, PyTorch, NumPyro, Stan,
bnlearn) also run at Step 12 --- with per-kind reach recorded as
explicit statuses (tbl.~4) --- so that a
model written once in this notation can be parsed, visualized, and
executed across the 25-step pipeline (0--24) without restating its
mathematics (Smékal and Friedman 2023; Heins et al. 2022; Felice et al.
2021).
\end{frame}

\end{document}
